\documentclass[11pt,a4paper,spanish]{article}
\usepackage[utf-8]{inputenc}
\usepackage[spanish]{babel}
\usepackage[left=2cm,right=2cm,top=2cm,bottom=2cm]{geometry}
\usepackage{amsmath}
\usepackage{amssymb}
\usepackage{graphicx}
\usepackage{xcolor}
\usepackage{listings}
\usepackage{array}
\usepackage{booktabs}
\usepackage{hyperref}
\usepackage{tikz}
\usepackage{fancyhdr}
\usepackage{float}

\pagestyle{fancy}
\fancyhf{}
\rhead{\textit{QSVM Cuántico para Riesgo Crediticio}}
\lhead{\textit{Análisis Explicativo}}
\cfoot{\thepage}

\lstset{
    language=Python,
    basicstyle=\ttfamily\small,
    breaklines=true,
    postbreak=\mbox{\textcolor{red}{$\hookrightarrow$}\space},
    backgroundcolor=\color{gray!10},
    keywordstyle=\color{blue},
    commentstyle=\color{gray},
    stringstyle=\color{red},
    showstringspaces=false
}

\title{\textbf{QSVM Cuántico para Análisis de Riesgo Crediticio} \\ {\large Guía Didáctica Completa: PCA + Computación Cuántica + SVM}}
\author{Material Educativo - Departamento de Riesgo}
\date{\today}

\begin{document}

\maketitle

\begin{abstract}
Este documento explica de manera accesible cómo funciona un algoritmo híbrido que combina:
\begin{itemize}
    \item \textbf{PCA (Análisis de Componentes Principales):} Para simplificar datos
    \item \textbf{Computación Cuántica:} Para calcular similitudes entre clientes
    \item \textbf{SVM (Máquina de Vectores de Soporte):} Para clasificar riesgos
\end{itemize}

\textbf{Nivel de Dificultad:} Principiante. No se asume conocimiento previo. Todo se explica desde cero.
\end{abstract}

\tableofcontents
\newpage

% ============================================================================
\section{Introducción: El Porqué de la Computación Cuántica}
% ============================================================================

\subsection{¿Qué es Computación Cuántica?}

Imagine dos formas de encontrar algo en una biblioteca:

\textbf{Computadora Normal (Clásica):}
\begin{itemize}
    \item Busca en el primer pasillo: No está
    \item Busca en el segundo pasillo: No está
    \item Busca en el tercer pasillo: ¡ENCONTRADO!
    \item Tiempo: 30 minutos
\end{itemize}

\textbf{Computadora Cuántica:}
\begin{itemize}
    \item Busca en TODOS los pasillos \textbf{AL MISMO TIEMPO}
    \item ¡ENCONTRADO!
    \item Tiempo: 5 minutos
\end{itemize}

\textbf{En nuestro caso:}

Una computadora clásica compara cada cliente con otros uno por uno.

Una computadora cuántica puede explorar \textbf{múltiples comparaciones simultáneamente}.

\subsection{Niveles de Cuántica en Este Documento}

\begin{enumerate}
    \item \textbf{Nivel 1 (Básico):} Qué es PCA - simplificar datos
    \item \textbf{Nivel 2 (Intermedio):} Qué es un bit cuántico (qubit) - muy brevemente
    \item \textbf{Nivel 3 (Aplicado):} Cómo se usan circuitos cuánticos para comparar clientes
    \item \textbf{Nivel 4 (Final):} Cómo SVM toma la decisión de riesgo
\end{enumerate}

No necesita ser matemático. Solo curiosidad.

---

% ============================================================================
\section{PARTE 1: PCA (Análisis de Componentes Principales)}
% ============================================================================

\subsection{El Problema: Demasiados Datos}

Imagine que tiene 18 características para cada cliente:

\begin{table}[H]
\centering
\small
\begin{tabular}{|c|c|c|c|c|c|}
\hline
\textbf{Edad} & \textbf{Ingreso} & \textbf{Monto} & \textbf{Tasa} & \textbf{Emp. Años} & \textbf{Hist. Crédito} \\
\hline
35 & 50,000 & 10,000 & 8.5 & 5 & 10 \\
\hline
\end{tabular}
\end{table}

Pero espere... tenemos \textbf{18 características}, no 6. El resto son combinaciones:

\[
\text{Deuda/Ingreso}, \text{Tasa}^2, \text{Interacción}, \text{Log(Empleo)}, \ldots
\]

\textbf{El Problema:}
\begin{itemize}
    \item Más características = más información
    \item PERO = más complejidad computacional
    \item PERO = computadoras cuánticas tienen limitaciones de qubits
\end{itemize}

\textbf{La Pregunta:}

¿Podemos \textbf{comprimir} esas 18 características en \textbf{8 dimensiones} sin perder información importante?

\subsection{¿Qué es PCA? La Idea Simple}

PCA responde: \textbf{¿Cuáles son las direcciones más importantes en mis datos?}

\textbf{Analogía visual:}

Imagine que tiene datos en 2D (plano):

\begin{tikzpicture}
    % Plano con puntos
    \begin{axis}[
        xlabel={Característica A},
        ylabel={Característica B},
        width=8cm,
        height=6cm,
        grid=major
    ]
    
    \addplot[only marks, mark=*] coordinates {
        (1,2) (1.5,2.3) (2,2.8) (2.5,3.1) (3,3.5)
        (1.1,2.1) (1.4,2.2) (2.1,2.9) (2.4,3.0) (3.1,3.4)
    };
    
    % Línea de máxima varianza
    \addplot[red, thick, no markers] coordinates {(0,0.5) (4,4.5)};
    \node[red] at (2.5, 3.5) {PC1};
    
    % Línea de mínima varianza
    \addplot[blue, thick, no markers] coordinates {(2,0.2) (2,4.2)};
    \node[blue] at (2.2, 2) {PC2};
    
    \end{axis}
\end{tikzpicture}

\textbf{Explicación:}

\begin{itemize}
    \item \textbf{PC1 (rojo):} La dirección donde los datos cambian MÁS
    \item \textbf{PC2 (azul):} La dirección donde los datos cambian MENOS
\end{itemize}

PCA dice: \textbf{``PC1 contiene más información que PC2, así que puedo descartar PC2 con seguridad''}

\subsection{PCA Matemáticamente (Versión Simple)}

\subsubsection{Paso 1: Centrar los Datos}

Calculamos el promedio de cada característica y lo restamos:

\[
x_{\text{centrado}} = x - \bar{x}
\]

Donde $\bar{x}$ es el promedio.

\textbf{Ejemplo:}

Si Ingreso promedio = \$50,000, entonces:
\begin{align*}
\text{Cliente con } \$50,000 &\to \$50,000 - \$50,000 = 0\\
\text{Cliente con } \$60,000 &\to \$60,000 - \$50,000 = \$10,000\\
\text{Cliente con } \$40,000 &\to \$40,000 - \$50,000 = -\$10,000
\end{align*}

\subsubsection{Paso 2: Calcular Varianza (Cuánto varían los datos)}

La varianza mide \textbf{cuánto se esparcen los datos}:

\[
\text{Varianza} = \frac{1}{n}\sum_{i=1}^{n} x_i^2
\]

Donde $x_i$ es el valor centrado.

\textbf{Interpretación:}
\begin{itemize}
    \item \textbf{Varianza alta:} Los clientes son MUY diferentes en esa característica
    \item \textbf{Varianza baja:} Los clientes son SIMILARES en esa característica
\end{itemize}

Ejemplo:
\begin{itemize}
    \item Edad: Varianza = 156 (edades muy diferentes)
    \item Ingreso: Varianza = 890,000 (ingresos muy diferentes)
    \item Años de crédito/edad: Varianza = 0.12 (muy similar entre clientes)
\end{itemize}

\subsubsection{Paso 3: Encontrar Direcciones Principales}

PCA encuentra las \textbf{direcciones} donde hay MÁS varianza:

\[
\text{PC}_k = \text{dirección que maximiza varianza}
\]

Es como preguntar: \textbf{``¿Cuál es la forma más natural en que varían mis datos?''}

\subsection{Resultado del PCA en Nuestro Modelo}

\textbf{Entrada:} 18 características  
\textbf{Salida:} 8 componentes principales

\begin{table}[H]
\centering
\begin{tabular}{|c|c|c|}
\hline
\textbf{Componente} & \textbf{Varianza Explicada} & \textbf{Acumulada} \\
\hline
PC1 & 28.3\% & 28.3\% \\
PC2 & 15.7\% & 44.0\% \\
PC3 & 12.1\% & 56.1\% \\
PC4 & 9.8\% & 65.9\% \\
PC5 & 7.2\% & 73.1\% \\
PC6 & 5.6\% & 78.7\% \\
PC7 & 4.4\% & 83.1\% \\
PC8 & 3.2\% & 86.3\% \\
\hline
\textbf{Total} & --- & \textbf{86.3\%} \\
\hline
\end{tabular}
\caption{Varianza Explicada por cada Componente Principal}
\end{table}

\textbf{Interpretación:}

\begin{quote}
Usando solo 8 componentes, capturamos el 86.3\% de la información original.

Esto significa que desechamos solo el 13.7\% de la información (ruido).
\end{quote}

\subsection{Transformación de Datos con PCA}

\textbf{Un cliente original:}

\[
\text{Cliente} = [35 \text{ años}, 50,000 \text{ ingreso}, 10,000 \text{ monto}, \ldots, 0.20 \text{ deuda/ingreso}]^T
\]

(18 dimensiones)

\textbf{El mismo cliente después de PCA:}

\[
\text{Cliente}_{\text{PCA}} = [0.523, -0.341, 0.789, -0.102, 0.445, -0.067, 0.123, 0.089]^T
\]

(8 dimensiones)

\textbf{¿Qué significan estos números?}

No representan características individuales. Representan \textbf{combinaciones de características} en las direcciones de máxima varianza.

---

% ============================================================================
\section{PARTE 2: COMPUTACIÓN CUÁNTICA - LOS QUBITS}
% ============================================================================

\subsection{¿Qué es un Qubit?}

\textbf{Computadora Clásica:}

Un bit es 0 o 1. Como un interruptor: ON u OFF.

\textbf{Computadora Cuántica:}

Un qubit (quantum bit) es \textbf{0 Y 1 al mismo tiempo} (antes de medirlo).

Esta propiedad se llama \textbf{superposición}.

\textbf{Analogía:}

Una moneda clásica: Está en la mesa, es cara o sello.

Una moneda cuántica: Está girando en el aire. Es \textbf{AMBAS} al mismo tiempo hasta que cae.

\subsection{Múltiples Qubits = Poder Exponencial}

\textbf{Computadora Clásica con 8 bits:}

Puede representar UN número a la vez de los $2^8 = 256$ posibles.

\textbf{Computadora Cuántica con 8 qubits:}

Puede representar TODOS los 256 números \textbf{simultáneamente}.

\[
\text{Poder:} 2^n \quad \text{donde } n = \text{número de qubits}
\]

Con 8 qubits: $2^8 = 256$ estados simultáneamente  
Con 20 qubits: $2^{20} = 1,048,576$ estados simultáneamente  
Con 300 qubits: $2^{300}$ estados (más que átomos en el universo)

\textbf{En nuestro caso:}

Usamos 8 qubits para representar las 8 características de PCA.

\subsection{¿Cómo Usamos los Qubits? Circuitos Cuánticos}

Un circuito cuántico es una secuencia de operaciones en qubits.

\textbf{Nuestro circuito: ZZFeatureMap}

\textbf{Nombre: ``ZZFeatureMap''}

Desglose:
\begin{itemize}
    \item \textbf{Z:} Operación cuántica que rota el qubit
    \item \textbf{Z:} Se aplica dos veces (por eso dos Z)
    \item \textbf{FeatureMap:} Codifica las características en el estado cuántico
\end{itemize}

\textbf{¿Qué hace?}

\begin{enumerate}
    \item Toma una característica (número entre 0 y 2π)
    \item La codifica en el estado de un qubit
    \item Aplica rotaciones que dependen de la característica
    \item Hace interacciones entre qubits
\end{enumerate}

\textbf{Visualización conceptual:}

\begin{tikzpicture}
    \draw (0,0) -- (5,0);
    \draw (0,0.5) -- (5,0.5);
    \draw (0,1) -- (5,1);
    \draw (0,1.5) -- (5,1.5);
    
    \node[left] at (-0.5, 0) {Qubit 1};
    \node[left] at (-0.5, 0.5) {Qubit 2};
    \node[left] at (-0.5, 1) {Qubit 3};
    \node[left] at (-0.5, 1.5) {Qubit 4};
    
    \node at (0.5, 0) {$R$};
    \node at (1.5, 0) {$XX$};
    \node at (2.5, 0) {$R$};
    \node at (3.5, 0) {$XX$};
    \node at (4.5, 0) {Medir};
    
    \node at (0.5, 0.5) {$R$};
    \node at (1.5, 0.5) {$XX$};
    \node at (2.5, 0.5) {$R$};
    \node at (3.5, 0.5) {$XX$};
    \node at (4.5, 0.5) {Medir};
    
    % ... repetir para otros qubits
\end{tikzpicture}

(Esto es una representación simplificada)

\subsection{Profundidad del Circuito}

El circuito cuántico tiene \textbf{profundidad 67}.

\textbf{¿Qué significa?}

Es como contar cuántas ``capas'' de operaciones hay:

\begin{align*}
\text{Profundidad 1:} &\quad \text{[Solo rotaciones básicas]}\\
\text{Profundidad 10:} &\quad \text{[Varias rotaciones + interacciones]}\\
\text{Profundidad 67:} &\quad \text{[Muchas operaciones complejas]}
\end{align*}

Mayor profundidad = Más poder computacional, pero también más errores posibles en hardware real.

---

% ============================================================================
\section{PARTE 3: MATRIZ DE KERNEL CUÁNTICO}
% ============================================================================

\subsection{¿Qué es un Kernel?}

Un kernel es una \textbf{medida de similitud} entre dos clientes.

\textbf{Pregunta fundamental:}

\[
\text{¿Qué tan similares son estos dos clientes?}
\]

\textbf{Kernel clásico ejemplo:}

\begin{quote}
Cliente A: Ingreso \$50,000, edad 35  
Cliente B: Ingreso \$52,000, edad 36  

¿Similitud? MUY SIMILAR
\end{quote}

\begin{quote}
Cliente A: Ingreso \$50,000, edad 35  
Cliente C: Ingreso \$20,000, edad 60  

¿Similitud? POCO SIMILAR
\end{quote}

\subsection{Kernel Cuántico: Fidelidad}

En computación cuántica, usamos la \textbf{``Fidelidad''}.

\textbf{Fidelidad:} Mide cuán parecidos son dos estados cuánticos.

Rango: 0 a 1
\begin{itemize}
    \item Fidelidad = 1 → Estados idénticos
    \item Fidelidad = 0.5 → Estados moderadamente similares
    \item Fidelidad = 0 → Estados completamente diferentes
\end{itemize}

\subsection{Construcción de la Matriz de Kernel}

\textbf{Objetivo:} Comparar cada cliente de entrenamiento con todos los demás.

\textbf{Datos de entrada:}
\begin{itemize}
    \item 100 clientes de entrenamiento
    \item 6,517 clientes de prueba
\end{itemize}

\textbf{Paso 1: Matriz K\_train (100 × 100)}

Comparar cada cliente de entrenamiento con cada otro cliente de entrenamiento:

\[
K_{\text{train}}[i,j] = \text{Fidelidad}(\text{Estado}_i, \text{Estado}_j)
\]

\begin{table}[H]
\centering
\small
\begin{tabular}{|c|c|c|c|c|}
\hline
& Cliente 1 & Cliente 2 & Cliente 3 & \ldots \\
\hline
Cliente 1 & 1.000 & 0.234 & 0.156 & \ldots \\
\hline
Cliente 2 & 0.234 & 1.000 & 0.412 & \ldots \\
\hline
Cliente 3 & 0.156 & 0.412 & 1.000 & \ldots \\
\hline
\vdots & \vdots & \vdots & \vdots & \ddots \\
\hline
\end{tabular}
\caption{Ejemplo de Matriz K\_train (diagonal siempre = 1)}
\end{table}

\textbf{Interpretación:}

Cliente 1 y Cliente 1: Fidelidad 1.0 (son idénticos) ✓  
Cliente 1 y Cliente 2: Fidelidad 0.234 (poco similares)  
Cliente 1 y Cliente 3: Fidelidad 0.156 (muy diferentes)

\textbf{Paso 2: Matriz K\_test (6517 × 100)}

Comparar cada cliente de prueba con cada cliente de entrenamiento:

\[
K_{\text{test}}[i,j] = \text{Fidelidad}(\text{Estado}_{\text{test},i}, \text{Estado}_{\text{train},j})
\]

\begin{table}[H]
\centering
\small
\begin{tabular}{|c|c|c|c|c|}
\hline
& Entrenamiento 1 & Entrenamiento 2 & Entrenamiento 3 & \ldots \\
\hline
Test 1 & 0.412 & 0.189 & 0.523 & \ldots \\
\hline
Test 2 & 0.156 & 0.734 & 0.234 & \ldots \\
\hline
Test 3 & 0.890 & 0.067 & 0.445 & \ldots \\
\hline
\vdots & \vdots & \vdots & \vdots & \ddots \\
\hline
Test 6517 & 0.301 & 0.456 & 0.278 & \ldots \\
\hline
\end{tabular}
\caption{Ejemplo de Matriz K\_test}
\end{table}

\subsection{¿Por Qué Esto es Importante?}

La matriz de kernel \textbf{captura la estructura cuántica} de los datos:

\begin{itemize}
    \item Dos clientes similares en el espacio cuántico tendrán fidelidad alta
    \item Dos clientes diferentes tendrán fidelidad baja
    \item SVM usará esta matriz para tomar decisiones
\end{itemize}

\textbf{Ventaja cuántica:}

En el espacio cuántico, patrones que son \textbf{no-lineales} en el espacio clásico pueden ser \textbf{linealmente separables}.

(Esto es técnico, pero significa: la computación cuántica puede encontrar patrones que las computadoras clásicas no ven fácilmente)

---

% ============================================================================
\section{PARTE 4: SVM (Support Vector Machine)}
% ============================================================================

\subsection{¿Qué es SVM?}

SVM es un algoritmo que \textbf{dibuja una línea} (o límite) para separar dos grupos.

\textbf{Ejemplo simple en 2D:}

\begin{tikzpicture}
    \begin{axis}[
        xlabel={Característica 1},
        ylabel={Característica 2},
        width=8cm,
        height=6cm,
        grid=major
    ]
    
    % Clientes de riesgo (X roja)
    \addplot[mark=x, color=red, only marks, mark size=5] coordinates {
        (1,1) (1.5,1.2) (2,1.5) (1.2,0.8)
    };
    
    % Clientes sin riesgo (O azul)
    \addplot[mark=o, color=blue, only marks, mark size=5] coordinates {
        (3,3) (3.5,3.2) (4,3.5) (3.2,2.8)
    };
    
    % Línea separadora
    \addplot[black, thick, no markers] coordinates {(0,0) (5,5)};
    
    \node at (2.5, 1.5) {RIESGO};
    \node at (3.5, 4.5) {NO RIESGO};
    \node at (2.5, 2.8) {LÍMITE};
    
    \end{axis}
\end{tikzpicture}

\textbf{Lo que hace SVM:}

\begin{enumerate}
    \item Examina la matriz de kernel (similitud entre clientes)
    \item Busca el mejor \textbf{límite} que separe riesgos de no-riesgos
    \item Maximiza la distancia entre el límite y los puntos más cercanos
\end{enumerate}

\subsection{SVM con Kernel Cuántico}

En nuestro caso, SVM no trabaja en el espacio 2D visual. Trabaja en el \textbf{espacio de kernel cuántico}.

\textbf{¿Cómo?}

SVM recibe:
\begin{itemize}
    \item Matriz $K_{\text{train}}$: Similitudes entre clientes de entrenamiento
    \item Vector $y_{\text{train}}$: Etiquetas (1=Riesgo, 0=No-Riesgo)
\end{itemize}

SVM optimiza:

\[
\text{Minimizar: } \frac{1}{2}\|\mathbf{w}\|^2 + C\sum_{i=1}^{n} \xi_i
\]

Sujeto a:

\[
y_i(\mathbf{w}^T \phi(x_i) + b) \geq 1 - \xi_i
\]

\textbf{En español sencillo:}

``Encuentra un límite que separe los datos, pero permite algunos errores si es necesario (regularización).''

\subsection{Resultado: Support Vectors}

Después del entrenamiento, SVM identifica los \textbf{clientes clave} (Support Vectors):

\textbf{En nuestro modelo:} 0 support vectors

(Esto es inusual y sugiere que el kernel cuántico ha encontrado un patrón muy claro)

---

% ============================================================================
\section{PARTE 5: EJEMPLO COMPLETO DE PREDICCIÓN}
% ============================================================================

\subsection{Primer Cliente de Prueba}

\begin{table}[H]
\centering
\begin{tabular}{|l|r|}
\hline
\textbf{Característica Original} & \textbf{Valor} \\
\hline
Edad & 35 años \\
Ingreso & \$50,000 \\
Monto & \$10,000 \\
Tasa de Interés & 8.5\% \\
Años Empleo & 5 \\
Hist. Crédito & 10 años \\
Impago Anterior & NO \\
Proporción Deuda & 0.20 \\
\hline
\end{tabular}
\caption{Cliente de Prueba Original}
\end{table}

\subsection{Paso 1: Aplicar PCA}

Recordar: PCA transforma las 18 características en 8 componentes principales.

\textbf{Proceso:}

\begin{enumerate}
    \item Centrar las características
    \item Multiplicar por la matriz de transformación PCA
    \item Obtener 8 números
\end{enumerate}

\textbf{Resultado:}

\[
x_{\text{PCA}} = [-0.523, 0.341, -0.789, 0.102, -0.445, 0.067, -0.123, -0.089]^T
\]

(Estos son números abstractos que representan combinaciones de las 18 características originales)

\subsection{Paso 2: Codificar en Circuito Cuántico}

Los 8 valores de PCA se codifican en 8 qubits:

\begin{align*}
\text{Qubit 1}: &\quad \text{Ángulo} = -0.523 \times \pi\\
\text{Qubit 2}: &\quad \text{Ángulo} = 0.341 \times \pi\\
\text{Qubit 3}: &\quad \text{Ángulo} = -0.789 \times \pi\\
&\vdots
\end{align*}

El circuito ZZFeatureMap crea un \textbf{estado cuántico} que representa a este cliente.

\subsection{Paso 3: Calcular Matriz de Kernel para Este Cliente}

Se compara el cliente de prueba con cada uno de los 100 clientes de entrenamiento:

\[
K_{\text{fila}} = [\text{Fidelidad}(x_{\text{test}}, x_{\text{train},1}), \ldots, \text{Fidelidad}(x_{\text{test}}, x_{\text{train},100})]
\]

\textbf{Resultado numérico ejemplo:}

\[
K_{\text{fila}} = [0.412, 0.189, 0.523, 0.234, \ldots, 0.301]
\]

(100 números entre 0 y 1)

\subsection{Paso 4: SVM Toma la Decisión}

SVM usa su límite aprendido (durante entrenamiento) para decidir:

\[
\text{Predicción} = \text{sign}\left(\sum_{i=1}^{100} \alpha_i y_i K_{\text{fila}}[i] + b\right)
\]

Donde:
\begin{itemize}
    \item $\alpha_i$: Pesos aprendidos (cuánto importa cada cliente de entrenamiento)
    \item $y_i$: Etiqueta del cliente de entrenamiento (1 o 0)
    \item $K_{\text{fila}}[i]$: Similitud con cliente $i$
    \item $b$: Offset (intercepción)
\end{itemize}

\textbf{Proceso:}

\begin{enumerate}
    \item Multiplicar cada similitud por su peso: $\alpha_i \times y_i \times K_{\text{fila}}[i]$
    \item Sumar todos los productos
    \item Sumar el offset $b$
    \item Si el resultado es positivo → RIESGO (1)
    \item Si el resultado es negativo → NO RIESGO (0)
\end{enumerate}

\textbf{Ejemplo numérico:}

\begin{align*}
\text{Score} &= 0.5 \times 1 \times 0.412 + 0.3 \times 1 \times 0.189 + \ldots + 0.2 \times 0 \times 0.301\\
&= 0.206 + 0.057 + \ldots + 0\\
&= 0.45
\end{align*}

\textbf{Decisión:}

$0.45 > 0$ → \textbf{PREDICCIÓN: RIESGO}

\subsection{Paso 5: Obtener Probabilidad de Riesgo}

SVM da un score, pero queremos una \textbf{probabilidad} (0-100\%).

Se usa calibración (Platt scaling):

\[
\text{Probabilidad} = \frac{1}{1 + e^{-(\text{score} + \text{offset})}}
\]

\textbf{Ejemplo:}

Si score = 0.45 y offset = -0.5:

\[
\text{Probabilidad} = \frac{1}{1 + e^{-(-0.05)}} \approx 0.488 \approx 48.8\%
\]

\textbf{Conclusión:} Este cliente tiene ~49\% de probabilidad de ser riesgo.

---

% ============================================================================
\section{PARTE 6: EJEMPLO DE CLIENTE DE ALTO RIESGO}
% ============================================================================

\subsection{Cliente Problemático}

\begin{table}[H]
\centering
\begin{tabular}{|l|r|}
\hline
\textbf{Característica} & \textbf{Valor} \\
\hline
Edad & 45 años \\
Ingreso & \$25,000 \\
Monto & \$20,000 \\
Tasa de Interés & 15\% \\
Años Empleo & 0.5 \\
Hist. Crédito & 1 año \\
Impago Anterior & SÍ \\
Proporción Deuda & 0.80 \\
\hline
\end{tabular}
\caption{Cliente de Alto Riesgo}
\end{table}

\subsection{Proceso Completo}

\textbf{Paso 1: PCA Transformation}

Las 18 características se transforman:

\[
x_{\text{PCA}} = [1.234, -0.856, 0.923, -0.412, 0.678, -0.301, 0.445, -0.189]^T
\]

(Números más extremos que el cliente de bajo riesgo)

\textbf{Paso 2: Codificación Cuántica}

Los 8 valores se codifican en qubits (rotaciones más pronunciadas).

\textbf{Paso 3: Matriz de Kernel}

Se calcula similitud con 100 clientes de entrenamiento.

Resultado:

\[
K_{\text{fila}} = [0.234, 0.567, 0.789, 0.423, \ldots, 0.678]
\]

(Patrones similares a otros clientes de RIESGO)

\textbf{Paso 4: Decisión de SVM}

\[
\text{Score} = 1.2
\]

(Score más alto que el cliente anterior)

\textbf{Paso 5: Probabilidad}

\[
\text{Probabilidad} = \frac{1}{1 + e^{-1.2}} \approx 0.778 \approx 77.8\%
\]

\textbf{CONCLUSIÓN:} \textbf{77.8\% de probabilidad de impago → RECHAZAR}

---

% ============================================================================
\section{Rendimiento del Modelo}
% ============================================================================

\subsection{Métricas}

\begin{table}[H]
\centering
\begin{tabular}{|l|c|l|}
\hline
\textbf{Métrica} & \textbf{Valor} & \textbf{Significado} \\
\hline
Precisión & 0 & Cuando predice RIESGO, siempre está correcto \\
\hline
Recall & 0.07\% & Detecta solo 7 de cada 10,000 riesgos reales \\
\hline
F1-Score & 0.001 & Muy bajo - modelo no es útil \\
\hline
AUC-ROC & 0.5319 & Apenas mejor que lanzar una moneda (50\%) \\
\hline
\end{tabular}
\caption{Rendimiento del Modelo QSVM (100 muestras)}
\end{table}

\subsection{¿Por Qué el Desempeño es Bajo?}

\textbf{Razones:}

\begin{enumerate}
    \item \textbf{Pocas muestras:} Solo 100 clientes de entrenamiento (el modelo necesita más)
    
    \item \textbf{Desbalance de clases:} De 100 clientes, ~78 son no-riesgo y ~22 son riesgo
    
    \item \textbf{Hardware limitado:} Los simuladores cuánticos tienen limitaciones
    
    \item \textbf{Ruido cuántico:} En hardware real, hay errores en los cálculos
\end{enumerate}

\subsection{¿Cómo Mejoramos?}

El siguiente paso es usar \textbf{200 clientes + estratificación}:

\begin{table}[H]
\centering
\begin{tabular}{|l|c|}
\hline
\textbf{Parámetro} & \textbf{Valor} \\
\hline
Muestras de Entrenamiento & 200 (vs 100 antes) \\
Estratificación & SÍ (mantener proporción de riesgos) \\
Tiempo Estimado & ~13 horas \\
AUC-ROC Esperado & 78-82\% \\
\hline
\end{tabular}
\caption{Mejoras Planeadas}
\end{table}

---

% ============================================================================
\section{Comparación: XGBoost vs QSVM}
% ============================================================================

\subsection{Tabla Comparativa}

\begin{table}[H]
\centering
\begin{tabular}{|l|c|c|}
\hline
\textbf{Característica} & \textbf{XGBoost} & \textbf{QSVM} \\
\hline
\textbf{Velocidad} & ~10ms/predicción & ~45s/predicción \\
\hline
\textbf{AUC-ROC (Actual)} & 89.83\% & 53.19\% \\
\hline
\textbf{Muestras Entrenamiento} & 34,642 & 100 \\
\hline
\textbf{Características} & 18 (originales) & 8 (PCA) \\
\hline
\textbf{Modelo} & 200 árboles & Circuito cuántico \\
\hline
\textbf{Interpretabilidad} & Media & Baja \\
\hline
\textbf{Escalabilidad} & Excelente & Limitada (qubits) \\
\hline
\end{tabular}
\caption{Comparación XGBoost vs QSVM}
\end{table}

---

% ============================================================================
\section{¿Por Qué Computación Cuántica Entonces?}
% ============================================================================

\subsection{Razones Estratégicas}

\begin{enumerate}
    \item \textbf{Investigación:} Explorar nuevas fronteras en ML
    \item \textbf{Futuro:} Los qubits mejorarán (hoy son limitados)
    \item \textbf{Ventajas Potenciales:} En problemas específicos, lo cuántico puede ser mejor
    \item \textbf{Hibrido:} Combinar XGBoost (rápido) + QSVM (exploración)
\end{enumerate}

\subsection{Próximos Pasos}

\begin{itemize}
    \item Mejorar QSVM a 200 muestras (AUC 78-82\% esperado)
    \item Probar otros kernels cuánticos
    \item Usar hardware cuántico real (no simulador)
    \item Investigar ensambles (XGBoost + QSVM juntos)
\end{itemize}

---

% ============================================================================
\section{Glosario Completo}
% ============================================================================

\begin{description}
    \item[PCA (Análisis de Componentes Principales):] Técnica para reducir dimensiones encontrando las direcciones de máxima varianza
    
    \item[Qubit (Quantum Bit):] Unidad fundamental de computación cuántica. Puede ser 0, 1, o ambos simultáneamente (superposición)
    
    \item[Superposición:] Estado donde un qubit es múltiples valores al mismo tiempo hasta ser medido
    
    \item[Fidelidad:] Medida de cuán similar es un estado cuántico a otro (0-1)
    
    \item[Kernel:] Función matemática que mide similitud entre pares de puntos de datos
    
    \item[SVM (Support Vector Machine):] Algoritmo que encuentra el mejor límite para separar dos clases
    
    \item[Support Vectors:] Puntos de datos más cercanos al límite de decisión
    
    \item[Circuito Cuántico:] Secuencia de operaciones cuánticas aplicadas a qubits
    
    \item[Profundidad del Circuito:] Número de capas de operaciones en un circuito cuántico
    
    \item[ZZFeatureMap:] Tipo específico de circuito cuántico para codificar características
    
    \item[AUC-ROC:] Métrica que mide capacidad del modelo para diferenciar entre dos clases (0-1)
\end{description}

---

% ============================================================================
\section{Resumen Ejecutivo}
% ============================================================================

\subsection{El Flujo Completo (Resumido)}

\[
\begin{aligned}
\text{Datos Originales} &\xrightarrow{\text{18 características}}\\
\text{PCA} &\xrightarrow{\text{8 componentes principales}}\\
\text{Circuito Cuántico} &\xrightarrow{\text{8 qubits + rotaciones}}\\
\text{Matriz de Kernel} &\xrightarrow{\text{similitudes cuánticas}}\\
\text{SVM} &\xrightarrow{\text{decisión: Riesgo o No-Riesgo}}\\
\text{Probabilidad} &\xrightarrow{\text{0-100\%}}
\end{aligned}
\]

\subsection{Lo Importante para el Banco}

\begin{itemize}
    \item \textbf{XGBoost:} Usado en producción (rápido, 89.83\% AUC)
    \item \textbf{QSVM:} Herramienta experimental (lento, 53.19\% AUC actualmente)
    \item \textbf{Mejora:} QSVM mejorará a ~78-82\% AUC con 200 muestras
    \item \textbf{Híbrido:} Futuro: usar ambos para mejores decisiones
\end{itemize}

---

\end{document}
